\section*{第 \ref{ch:g8:speed} 章——正比例、速度与平均数}

\begin{solution}{exo:g8:speed:1}
$\frac{150}{2} = 75$ km/h. \quad
$45$ 分钟 $= 0.75$ 小时：$\frac{27}{0.75} = 36$ km/h。 \quad
$\frac{100}{12.5} = 8$ m/s.
\end{solution}

\begin{solution}{exo:g8:speed:2}
$2 \times 85 = 170$ km. \quad
$40$ 分钟 $= \frac23$ 小时：$90 \times \frac23 = 60$ km。 \quad
$\frac{315}{90} = 3.5$ 小时 $= 3$ 小时 $30$ 分钟。
\end{solution}

\begin{solution}{exo:g8:speed:3}
$72 \div 3.6 = 20$ m/s. \quad
$5 \times 3.6 = 18$ km/h. \quad
$108 \div 3.6 = 30$ m/s.
\end{solution}

\begin{solution}{exo:g8:speed:4}
$600 \div 15 = 40$ 分钟。 \quad
$2$ 小时 $30 = 150$ 分钟：$15 \times 150 = 2250$ L。
\end{solution}

\begin{solution}{exo:g8:speed:5}
$\frac{3.30}{1.2} = 2.75$ 欧元/千克，对比
$\frac{2.16}{0.8} = 2.70$ 欧元/千克：第二种报价
略微更划算。
\end{solution}

\begin{solution}{exo:g8:speed:6}
\[
\bar v = \frac{25 \times 12.4 + 35 \times 10}{60}
= \frac{310 + 350}{60} = \frac{660}{60} = 11 .
\]
\end{solution}

\begin{solution}{exo:g8:speed:7}
$d = 340 \times 6 = 2040$ m $\approx 2$ km.
\end{solution}

\begin{solution}{exo:g8:speed:8}
距离：$2 \times 5 = 10$ km，再 $1.5 \times 4 = 6$ km，总计
$16$ km。所用时间：$2 + 0.5 + 1.5 = 4$ 小时。\omterm{def:g8:speed:def}{平均速度}：
$\frac{16}{4} = 4$ km/h.
\end{solution}

\begin{solution}{exo:g8:speed:9}
\omterm{def:g8:speed:average}{加权平均数}：
$\frac{3 \times 15 + 2 \times 9 + 1 \times 12}{3 + 2 + 1}
= \frac{45 + 18 + 12}{6} = \frac{75}{6} = 12.5$.
不\omterm{def:g8:speed:average}{加权平均数}：$\frac{15 + 9 + 12}{3} = 12$。两者不同，是因为
权数使\omterm{def:g6:fractions:def}{分数} $15$ 的权重
是\omterm{def:g6:fractions:def}{分数} $12$ 的三倍。
\end{solution}

\begin{solution}{exo:g8:speed:10}
\emph{1.} 时间：$\frac{120}{80} = 1.5$ 小时和 $\frac{120}{120} = 1$
小时。\omterm{def:g8:speed:def}{平均速度}：$\frac{240}{2.5} = 96$ km/h。

\emph{2.} 汽车在 $80$ km/h 下花费\emph{更多时间}（$1.5$ 小时）
比在 $120$ km/h 下（$1$ 小时）多：较慢的速度在
按时间加权的平均数中权重更大，使其低于中点 $100$。
\end{solution}

\begin{solution}{exo:g8:speed:11}
两人之间的 $18$ km 间距以 $5 + 4 = 9$ km/h 的速度缩小：
$\frac{18}{9} = 2$ 小时后相遇，时间为 $12{:}00$。安娜此时已走过
$2 \times 5 = 10$ km：两人在距安娜的村庄 $10$ km 处相遇（
且距鲍里斯的村庄 $8$ km——检验：$10 + 8 = 18$）。
\end{solution}

\begin{solution}{pb:g8:speed:1}
\textbf{1.} 由\cref{def:g8:speed:def}可知，$t_1 = \frac{d}{v_1}$ 与
$t_2 = \frac{d}{v_2}$；全程所用时间为
$\frac{d}{v_1} + \frac{d}{v_2}$，路程为 $2d$。

\textbf{2.} 从\omterm{def:g4:fractions:def}{分母}中提出 $d$，并将其
约去：
\[
\bar v
= \frac{2d}{d\left(\frac{1}{v_1} + \frac{1}{v_2}\right)}
= \frac{2}{\ \frac{1}{v_1} + \frac{1}{v_2}\ } .
\]
括号内通分：
$\frac{1}{v_1} + \frac{1}{v_2} = \frac{v_2 + v_1}{v_1 v_2}$，并
除以一个数即乘以它的\omterm{def:g8:fractions:inverse}{倒数}
（\cref{thm:g8:fractions:div}）：
\[
\bar v = 2 \times \frac{v_1 v_2}{v_1 + v_2}
= \frac{2\, v_1 v_2}{v_1 + v_2} .
\]

\textbf{3.} $\dfrac{2 \times 30 \times 15}{30 + 15}
= \dfrac{900}{45} = 20$ km/h，与
\cref{ex:g8:speed:twolegs}中的结果一致；
$\dfrac{2 \times 80 \times 120}{80 + 120} = \dfrac{19\,200}{200}
= 96$ km/h，与\cref{exo:g8:speed:10}中求得的结果一致。

\textbf{4.} \omterm{def:g8:speed:def}{平均速度}只取决于两个速度，与
$d$ 无关：$1$ km 的往返与 $100$ km 的往返，
在相同的两种速度下，\omterm{def:g8:speed:def}{平均速度}恰好相同。（这正是
公式能够在不同长度的行程上检验的原因。）

\textbf{5.} 若在每种速度下行驶时间均为 $T$，则路程分别为 $v_1 T$
和 $v_2 T$，于是
\[
\bar v = \frac{v_1 T + v_2 T}{2T}
= \frac{(v_1 + v_2)\,T}{2T}
= \frac{v_1 + v_2}{2} .
\]
速度的平均数总是按\emph{时间}加权
（\cref{def:g8:speed:average}）。时间相同时，两个速度
权重相等：即普通的中点。路程相等时，
较慢的速度占用\emph{更多}时间，因而权重更大——
平均数便向它偏移。

\textbf{6.} 对 $(30, 15)$，有 $H = 20$ 与 $A = 22.5$。对
$(80, 120)$，有 $H = 96$ 与 $A = 100$。
两次都是算术平均数更大。

\textbf{7.} 以公\omterm{def:g4:fractions:def}{分母} $2(v_1 + v_2)$ 通分：
\[
A - H
= \frac{(v_1 + v_2)^2 - 4\, v_1 v_2}{2\,(v_1 + v_2)} ,
\]
\omterm{def:g4:fractions:def}{分子}利用
\cref{pb:g8:equations:1}中的恒等式展开：
\[
(v_1 + v_2)^2 - 4 v_1 v_2
= v_1^2 + 2 v_1 v_2 + v_2^2 - 4 v_1 v_2
= v_1^2 - 2 v_1 v_2 + v_2^2
= (v_1 - v_2)^2 .
\]

\textbf{8.} 平方永远非负，且 $2(v_1 + v_2)$ 是
正数，故 $A - H \geq 0$，即 $H \leq A$，对每一对
正速度都成立。取等要求 $(v_1 - v_2)^2 = 0$，即
$v_1 = v_2$：只要两速度不同，往返的
\omterm{def:g8:speed:def}{平均速度}便严格低于中点——这一意外结论
本身便是一条定理。

\textbf{9.} 计算乘积：
\[
H \times A
= \frac{2\, v_1 v_2}{v_1 + v_2} \times \frac{v_1 + v_2}{2}
= v_1 v_2 ,
\]
约去因子 $2$ 与 $v_1 + v_2$ 后得出。

\textbf{10.} 代入公式：$\bar v = \dfrac{2 \times 40 \times 60}
{40 + 60} = \dfrac{4800}{100} = 48$ km/h。另一种算法：
$\frac{60}{40} = 1.5$ h 与 $\frac{60}{60} = 1$ h，故 $120$ km 用时
$2.5$ h：$\frac{120}{2.5} = 48$ km/h。两者一致。

\textbf{11.} 要以 $30$ km/h 的\omterm{def:g8:speed:def}{平均速度}跑完 $30$ km，意味着总时间
至多为 $\frac{30}{30} = 1$ h。前 $15$ km 以 $15$ km/h
的速度行进，已需 $\frac{15}{15} = 1$ h：\emph{全部}
时间预算已用尽，全程尚余\omterm{def:g3:division:half}{一半}。

\textbf{12.} 以 $90$ km/h 行驶后半程需
$\frac{15}{90} = \frac16$ h，故全程用时 $\frac76$ h，于是
\[
\bar v = \frac{30}{\ 7/6\ } = \frac{180}{7} \approx 25.7
\text{ km/h} < 30 .
\]
无论后半程速度为何，其时长都是正数，
故总时间超过 $1$ h，\omterm{def:g8:speed:def}{平均速度}仍
严格低于 $30$ km/h。（用公式：$\frac{2 \times 15
\times v_2}{15 + v_2} = 30$ 则须 $30 v_2 = 450 + 30 v_2$，
即 $0 = 450$——不存在满足要求的速度 $v_2$。）

\textbf{13.} \omterm{def:g8:speed:def}{平均速度} $20$ km/h 允许用时 $\frac{30}{20} = 1.5$ h；
第一小时后还剩 $15$ km，只有 $0.5$ h 可用：她需
以 $\frac{15}{0.5} = 30$ km/h 行驶。验算：
$H(15, 30) = \frac{2 \times 15 \times 30}{45} = \frac{900}{45}
= 20$ km/h.

\textbf{14.} 每分钟，两龙头分别注满水池的 $\frac{1}{20}$ 与
$\frac{1}{30}$，合计每分钟注入
\[
\frac{1}{20} + \frac{1}{30}
= \frac{3}{60} + \frac{2}{60}
= \frac{5}{60} = \frac{1}{12}
\]
水池\omterm{def:g2:measure:units}{容积}：即注满需 $12$ min。而
$H(20, 30) = \frac{2 \times 600}{50} = 24 = 2 \times 12$：两个水龙头
合用时恰为它们各自单独用时的 \emph{\omterm{pb:g8:speed:1}{调和平均数}的\omterm{def:g3:division:half}{一半}}，
速率相加，故仍属\omterm{pb:g8:speed:1}{调和平均数}的范畴。

\textbf{15.} 往返路程相等，故
\[
\bar v = \frac{2 \times 900 \times 700}{900 + 700}
= \frac{1\,260\,000}{1\,600} = 787.5 \text{ km/h},
\]
低于无风时的 $800$ km/h。一句话概括：飞机
逆风飞行所用时间\emph{更长}，于是较慢的那段
在按时间加权的平均数中权重更大——风
只会使往返耗时更长。
\end{solution}
